\documentclass[a4paper,12pt]{article}
\usepackage[utf8]{inputenc}
\usepackage[T1]{fontenc}
\usepackage[italian]{babel}
\usepackage[sfdefault]{atkinson}
\renewcommand*\familydefault{\sfdefault}
\usepackage{float}
\usepackage{microtype}
\usepackage{geometry}
\usepackage{setspace}
\usepackage{enumitem}
\usepackage{titlesec}
\usepackage{chngpage}
\usepackage{tocloft}
\usepackage{longtable}
\usepackage{array}
\usepackage{graphicx}
\usepackage{fancyhdr}
\usepackage[table]{xcolor}
\usepackage{color,soul}
\usepackage[most]{tcolorbox}
\usepackage{nameref}
\usepackage{changepage} 
\usepackage[colorlinks=true]{hyperref}

% --- Definizione Colori ---
\definecolor{zpusblack}{RGB}{0, 0, 0}

\hypersetup{
    linkcolor=black,
    urlcolor=blue
}

\definecolor{lightblack}{gray}{0.35}
\newcommand{\glossario}[1]{\textit{#1}\textsubscript{\textbf{\textit{\textcolor{lightblack}{G}}}}}
\newcommand{\ped}[1]{\textsubscript{#1}}

% --- Header e Footer ---
\pagestyle{fancy}
\setlength{\headwidth}{\textwidth}
\fancyhfoffset[L,R]{0pt}
\lhead{\rightmark}
\rhead{7-ZPUs}
\lfoot{Manuale Utente - DIPReader}
\rfoot{\thepage}
\cfoot{}
\renewcommand{\headrulewidth}{0.8pt}
\renewcommand{\footrulewidth}{0.8pt}

\renewcommand{\contentsname}{Indice}
\renewcommand{\listoftables}{Elenco delle tabelle}
\renewcommand{\listoffigures}{Elenco delle immagini}

\geometry{margin=2.5cm}
\setstretch{1.1}

% --- Configurazione Titoli (Nero e Fix Horizontal Mode) ---
\titleclass{\subsubsubsection}{straight}[\subsection]
\newcounter{subsubsubsection}[subsubsection]
\renewcommand\thesubsubsubsection{\thesubsubsection.\arabic{subsubsubsection}}

\titleformat{\section}{\LARGE\bfseries\textcolor{black}}{\thesection}{1em}{}
\titleformat{\subsection}{\Large\bfseries\textcolor{black}}{\thesubsection}{1em}{}
\titleformat{\subsubsection}{\large\bfseries\textcolor{black}}{\thesubsubsection}{1em}{}
\titleformat{\subsubsubsection}{\normalfont\normalsize\bfseries\textcolor{black}}{\thesubsubsubsection}{1em}{}

\titlespacing*{\section}{0pt}{3.5ex plus 1ex minus .2ex}{2.3ex plus .2ex}
\titlespacing*{\subsection}{0pt}{3.25ex plus 1ex minus .2ex}{1.5ex plus .2ex}

\makeatletter
\def\toclevel@subsubsubsection{4}
\def\l@subsubsubsection{\@dottedtocline{4}{7em}{4em}}
\makeatother

\setcounter{secnumdepth}{4}
\setcounter{tocdepth}{4}
\setlength\parindent{0pt}

% -----------------------------------------------------------------
% INIZIO DOCUMENTO
% -----------------------------------------------------------------
\begin{document}

\pagenumbering{arabic}
\setcounter{page}{1}

% --- Frontespizio ---
\begin{center}
    \includegraphics[width=9.5cm]{../../assets/logo7ZPUs.jpg}\\
    \small\hspace{10cm} 7zpus.swe@gmail.com\\
    \vspace{0.5cm}
    \LARGE \textbf{Manuale Utente}\\
    \vspace{0.3cm}
\end{center}

\vspace{0.3cm}
\hrule
\vspace{1cm}

% --- Tabella di Versionamento ---
\section*{Tabella di Versionamento}
\begin{table}[H]
    \begin{adjustwidth}{-1cm}{-1cm}
    \centering
    \begin{tabular}{|c|c|c|c|l|}
        \hline
        \textbf{Versione} & \textbf{Data} & \textbf{Autore}  & \textbf{Verificatore} & \textbf{Descrizione} \\
        \hline
        0.1.0 & 2026/04/09 & Georgescu Diana & Gingillino Aaron & \begin{tabular}[c]{@{}l@{}} Creazione del documento\\ e stesura iniziale \end{tabular} \\
        \hline
    \end{tabular}
    \end{adjustwidth}
\end{table}

\newpage
\tableofcontents

\newpage
\listoffigures
\listoftables

\newpage

% --- Sezione 1: Introduzione (Ripristinata) ---
\section{Introduzione}

\subsection{Obiettivo del documento}
L’obiettivo del presente documento è descrivere in modo rigoroso le funzionalità e l'utilizzo del prodotto \textit{DIPReader}.\\
L’analisi guiderà l'utente attraverso il contesto applicativo e le procedure operative per una gestione corretta del sistema.

\subsection{Glossario}
Ogni termine tecnico è definito nell'apposito documento: 
\href{https://7-zpus.github.io/Docs/assets/3_PB/DocumentiInterni/Glossario.pdf}{Glossario (v2.0.0)}.

\subsection{Riferimenti}

    \subsubsection{Riferimenti normativi}
    \begin{itemize}
        \item \href{https://github.com/7-ZPUs/Docs/blob/main/3_PB/DocumentiInterni/NormeDiProgetto.pdf}{Norme di Progetto v2.0.0}
        \item Standard SDD/IEEE/1016:2009
    \end{itemize}

    \subsubsection{Riferimenti informativi}
    \begin{itemize}
        \item \href{https://angular.io/}{Angular Documentation}
        \item \href{https://www.electronjs.org/}{Electron Documentation}
    \end{itemize}

\newpage

% --- Sezione 2: Tecnologie Utilizzate (Ripristinata) ---
\section{Tecnologie utilizzate}
Il sistema \textit{DIPReader} è sviluppato utilizzando tecnologie moderne per garantire prestazioni e sicurezza:
\begin{itemize}
    \item \textbf{Angular}: Framework per l'interfaccia utente basato su \textit{Signals}.
    \item \textbf{Electron}: Ambiente desktop per l'esecuzione cross-platform.
    \item \textbf{SQLite}: Motore di database locale per la persistenza dei dati.
\end{itemize}

\newpage

% --- Sezione 3: Contesto del Sistema (Ripristinata) ---
\section{Contesto del sistema}
\subsection{Modello architetturale}
Il software segue un modello di visualizzazione dei dati gerarchico.
\subsection{Descrizione degli attori esterni}
L'utente che consulta i pacchetti di conservazione.
\subsection{Descrizione dei sistemi esterni}
File system locale per l'archiviazione dei DIP.

\newpage

% --- Sezione 4: Guida all'Utilizzo (Ripristinata) ---
\section{Guida all'Utilizzo}
In questa sezione verranno dettagliate le operazioni di navigazione dell'albero DIP e la consultazione dei documenti tramite il visualizzatore PDF integrato.

\subsection{Esplorazione ad albero}
In qualsiasi momento è possibile vedere a sinistra una sezione dedicata alla vista ad albero del file DIP. Ogni riga rappresenta un raggruppamento oppure un file. Se una riga presenta una tabulazione, allora appartiene alla prima riga sopra che non presenta l'indentazione. La gerarchia dei raggruppamenti è la seguente:
\begin{itemize}
    \item Il nodo radice DIP contiene le classi documentali
    \item Le classi documenti contengono i processi associati (pallino rosso)
    \item I processi, che sono numerati in modo ordinato, contengono i file principali (pallino verde)
    \item I file principali contengono gli allegati (pallino azzurro) 
\end{itemize}
È possibile navigare questa sezione tramite il TAB. Se si desidera è possibile chiudere questa sezione tramite l'apposita linguetta "Albero" adiacente alla sezione. 

\subsection{Ricerca}
Per accedere a questa sezione è possibile premere il pulsante "Ricerca" in alto a destra, contrassegnato da una lente di ingrandimento. 
Questa sezione permette di ricercare all'interno del DIP tramite le seguenti modalità:
\begin{itemize}
    \item Per cercare all'interno del nome della classe o per ID del processo è possibile utilizzare la barra di ricerca, specificando tramite il pulsante d'opzione su quale campo si vuole compiere la ricerca. Per abilitare la ricerca semantica è possibile spuntare l'apposita casella, ora la ricerca includerà anche i risultati i cui campi specificati contengono termini nello stesso campo semantico, ovvero di significato simile, a quelli immessi nella barra di ricerca.
    \item A sinistra dei risultati è possibile aggiungere dei filtri pertinenti ai relativi metadati. Questi filtri possono essere nascosti tramite apposita linguetta, etichettata "filtri". 
\end{itemize}


\end{document}